\ifdefined\inMainDoc\else
\documentclass[12pt,a4paper]{article}
% ================================================================
%  PRÉAMBULE COMMUN - ANNALES CORRIGÉES
%  Université de Kara - FaST - Licence Fondamentale MPC
%  Design optimisé impression noir et blanc
% ================================================================

% -- ENCODAGE & LANGUE --
\usepackage[utf8]{inputenc}
\usepackage[T1]{fontenc}
\usepackage[french]{babel}
\usepackage{lmodern}
\usepackage{microtype}

% -- MISE EN PAGE --
\usepackage[a4paper, top=2.8cm, bottom=2.5cm, left=2.5cm, right=2.5cm, headheight=14pt]{geometry}
\usepackage{setspace}
\setstretch{1.2}
\usepackage{parskip}
\setlength{\parindent}{1.2em}
\setlength{\parskip}{4pt}

% -- MATHS --
\usepackage{amsmath, amssymb, mathtools, amsthm}
\usepackage{mathrsfs}
\usepackage{dsfont}
\usepackage{bm}
\usepackage{cancel}
\usepackage{textcomp}
% Définition de \oiint (intégrale de surface fermée) sans package supplémentaire
\newcommand{\oiint}{\mathop{\iint\!\!\!\!\!\!\!\bigcirc}\limits}

% -- PHYSIQUE & CHIMIE --
\usepackage{physics}
\usepackage{siunitx}
% Lorsque physics est chargé, siunitx omet \qty. On le restaure ici :
\AtBeginDocument{\RenewCommandCopy\qty\SI}
\sisetup{locale = FR, detect-all, output-decimal-marker={,}}
% Commande de substitution pour les unités posant problème avec pdflatex
% (ohm, degreeCelsius, micro, angstrom, percent utilisent des caractères non-ASCII)
\newcommand{\qtyohm}[1]{\ensuremath{\num{#1}\,\Omega}}
\newcommand{\qtydegc}[1]{\ensuremath{\num{#1}\,{}^\circ\text{C}}}
\newcommand{\qtyangstrom}[1]{\ensuremath{\num{#1}\,\text{\AA}}}
\newcommand{\qtypercent}[1]{\ensuremath{\num{#1}\,\%}}
\newcommand{\qtyum}[1]{\ensuremath{\num{#1}\,\mu\text{m}}}
\usepackage[version=4]{mhchem}

% -- GRAPHISMES --
\usepackage{graphicx}
\usepackage{float}
\usepackage{caption}
\usepackage{subcaption}
\usepackage{wrapfig}
\usepackage{tikz}
\usetikzlibrary{calc, arrows.meta, patterns, shapes.geometric}
\usepackage{pgfplots}
\pgfplotsset{compat=1.18}

% -- TABLEAUX --
\usepackage{booktabs}
\usepackage{array}
\usepackage{tabularx}
\usepackage{multirow}

% -- LISTES --
\usepackage{enumitem}
\setlist[enumerate]{topsep=3pt, itemsep=2pt, parsep=0pt}
\setlist[itemize]{topsep=3pt, itemsep=2pt, parsep=0pt, label={.}}

% -- BOÎTES (noir et blanc) --
\usepackage{mdframed}
\usepackage{tcolorbox}
\tcbuselibrary{skins, breakable, theorems}

% Boîte EXERCICE : encadré avec filet épais, fond blanc
\newmdenv[
  linewidth=1.2pt,
  linecolor=black,
  roundcorner=3pt,
  skipabove=10pt,
  skipbelow=6pt,
  innertopmargin=8pt,
  innerbottommargin=8pt,
  innerleftmargin=10pt,
  innerrightmargin=10pt,
  backgroundcolor=white,
  frametitle={\bfseries\small Exercice},
  frametitlealignment=\raggedright,
  frametitleaboveskip=4pt,
  frametitlebelowskip=4pt,
]{boiteexercice}

% Boîte CORRIGÉ : fond gris très clair + filet fin
\newmdenv[
  linewidth=0.6pt,
  linecolor=black,
  leftline=true,
  rightline=false,
  topline=false,
  bottomline=false,
  linecolor=black,
  innerleftmargin=12pt,
  innerrightmargin=8pt,
  innertopmargin=6pt,
  innerbottommargin=6pt,
  backgroundcolor=gray!8,
  skipabove=4pt,
  skipbelow=8pt,
]{boitecorrige}

% Boîte RÉSUMÉ DE COURS : double encadré
\newmdenv[
  linewidth=0.8pt,
  linecolor=black,
  roundcorner=2pt,
  backgroundcolor=gray!12,
  innertopmargin=10pt,
  innerbottommargin=10pt,
  innerleftmargin=12pt,
  innerrightmargin=12pt,
  skipabove=12pt,
  skipbelow=8pt,
]{boiteresume}

% Boîte ATTENTION/REMARQUE : barre gauche épaisse
\newmdenv[
  leftline=true,
  rightline=false,
  topline=false,
  bottomline=false,
  linewidth=3pt,
  linecolor=black,
  backgroundcolor=gray!5,
  innerleftmargin=12pt,
  innerrightmargin=8pt,
  innertopmargin=5pt,
  innerbottommargin=5pt,
  skipabove=6pt,
  skipbelow=6pt,
]{boiteremarque}

% -- DIFFICULTÉ DES EXERCICES --
\newcommand{\facile}{$\star$}
\newcommand{\moyen}{$\star\!\star$}
\newcommand{\difficile}{$\star\!\star\!\star$}
\newcommand{\niveau}[1]{\hfill\textbf{#1}\hspace{2pt}}

% -- EN-TÊTES ET PIEDS DE PAGE --
\usepackage{fancyhdr}
\pagestyle{fancy}
\fancyhf{}
\renewcommand{\headrulewidth}{0.5pt}
\renewcommand{\footrulewidth}{0.3pt}

% -- HYPERLIENS --
\usepackage[hidelinks]{hyperref}
\hypersetup{
    colorlinks=false,
    pdftitle={Annale Corrigée - Licence MPC - Université de Kara}
}

% -- TITRES --
\usepackage{titlesec}
\titleformat{\section}{\Large\bfseries}{\thesection.}{0.8em}{}[\vspace{2pt}\hrule\vspace{4pt}]
\titleformat{\subsection}{\large\bfseries}{\thesubsection.}{0.7em}{}
\titleformat{\subsubsection}{\normalsize\bfseries}{\thesubsubsection.}{0.6em}{}

% -- THÉORÈMES --
\theoremstyle{definition}
\newtheorem{defn}{Définition}[section]
\newtheorem{prop}{Propriété}[section]
\newtheorem{thm}{Théorème}[section]
\newtheorem{corol}{Corollaire}[section]
\newtheorem{exemple}{Exemple}[section]

% -- COMMANDES MATHÉMATIQUES UTILES --
\newcommand{\vect}[1]{\overrightarrow{#1}}
\newcommand{\R}{\mathbb{R}}
\newcommand{\N}{\mathbb{N}}
\newcommand{\Z}{\mathbb{Z}}
\newcommand{\C}{\mathbb{C}}
\renewcommand{\d}{\mathrm{d}}
\newcommand{\deriv}[2]{\dfrac{\d #1}{\d #2}}
\newcommand{\dpart}[2]{\dfrac{\partial #1}{\partial #2}}
\newcommand{\rot}{\overrightarrow{\mathrm{rot}}\,}
\renewcommand{\div}{\mathrm{div}\,}
\renewcommand{\grad}{\overrightarrow{\mathrm{grad}}\,}
\newcommand{\e}{\mathrm{e}}
\newcommand{\ii}{\mathrm{i}}
\renewcommand{\Tr}{\mathrm{Tr}}

% Numérotation des équations par section
\numberwithin{equation}{section}

% -- RÉSULTATS ENCADRÉS --
\newcommand{\res}[1]{\[\boxed{#1}\]}

% -- REMARQUE INLINE --
\newcommand{\rmq}[1]{\begin{boiteremarque}\textbf{Remarque.} #1\end{boiteremarque}}

% -- MISE EN VALEUR DU RÉSULTAT FINAL --
\newcommand{\reponse}[1]{%
  \begin{center}
    \fbox{\begin{minipage}{0.92\textwidth}\centering #1\end{minipage}}
  \end{center}%
}


\lhead{\small\textit{Programmation en langage C - Licence 1}}
\rhead{\small\textit{FaST, Université de Kara}}
\cfoot{\thepage}

% Pour le code source C
\usepackage{listings}
\lstset{
  language=C,
  basicstyle=\ttfamily\small,
  keywordstyle=\bfseries,
  commentstyle=\itshape,
  stringstyle=\ttfamily,
  numbers=left,
  numberstyle=\tiny,
  numbersep=8pt,
  frame=single,
  breaklines=true,
  tabsize=4,
  showstringspaces=false,
  literate={é}{{\'{e}}}1 {è}{{\`{e}}}1 {ê}{{\^{e}}}1 {ë}{{\"e}}1
           {à}{{\`a}}1 {â}{{\^a}}1 {ù}{{\`u}}1 {û}{{\^u}}1
           {î}{{\^i}}1 {ï}{{\"i}}1 {ô}{{\^o}}1 {\oe{}}{{\oe}}1
           {ç}{{\c{c}}}1
}

\begin{document}
\fi

\begin{center}
  {\LARGE\bfseries PROGRAMMATION EN LANGAGE C}\\[0.3cm]
  {\normalsize Licence 1 - Mathématiques, Physique et Chimie}\\[0.1cm]
  \rule{0.7\textwidth}{0.8pt}
\end{center}

\section*{Résumé du cours}

\begin{boiteresume}

\textbf{1. Structure d'un programme C.}
\begin{lstlisting}
#include <stdio.h>
int main() {
    /* instructions */
    return 0;
}
\end{lstlisting}

\textbf{2. Types de base.} \texttt{int} (entier), \texttt{float} (flottant simple précision), \texttt{double} (double précision), \texttt{char} (caractère). Déclaration : \texttt{type nom = valeur;}.

\textbf{3. Entrées/Sorties.} \texttt{printf("\%d \%f \%c \%s", i, x, c, s);} affiche entier, flottant, caractère, chaîne. \texttt{scanf("\%d", \&i);} lit un entier.

\textbf{4. Structures de contrôle.}
\begin{itemize}[nosep]
  \item Conditionnelle : \texttt{if (cond) \{...\} else \{...\}}
  \item Boucle \texttt{for} : \texttt{for (init; cond; increment) \{...\}}
  \item Boucle \texttt{while} : \texttt{while (cond) \{...\}}
  \item Boucle \texttt{do...while} : \texttt{do \{...\} while (cond);}
\end{itemize}

\textbf{5. Tableaux.} \texttt{int tab[10];} déclare un tableau de 10 entiers. Accès : \texttt{tab[i]} pour l'indice $i$ (de 0 à $n-1$).

\textbf{6. Fonctions.} \texttt{type\_retour nom(params) \{...; return val;\}}. Passage par valeur (copie) ou par pointeur (modification de l'original).

\textbf{7. Pointeurs.} \texttt{int *p = \&a;} : $p$ contient l'adresse de $a$. \texttt{*p} accède à la valeur pointée. \texttt{p++} passe à l'élément suivant.

\end{boiteresume}

\newpage
\section{Bases du langage C}

\subsection*{Exercice 1 - Variables et opérations \niveau{\facile}}

\begin{boiteexercice}
Écrire un programme C qui :
\begin{enumerate}
  \item Demande à l'utilisateur de saisir deux entiers $a$ et $b$.
  \item Calcule et affiche leur somme, différence, produit et quotient entier.
  \item Affiche également le reste de la division entière de $a$ par $b$.
\end{enumerate}
\end{boiteexercice}

\begin{boitecorrige}
\begin{lstlisting}
#include <stdio.h>

int main() {
    int a, b;

    printf("Entrez deux entiers a et b : ");
    scanf("%d %d", &a, &b);

    printf("Somme      : %d + %d = %d\n", a, b, a + b);
    printf("Difference : %d - %d = %d\n", a, b, a - b);
    printf("Produit    : %d * %d = %d\n", a, b, a * b);

    if (b != 0) {
        printf("Quotient   : %d / %d = %d\n", a, b, a / b);
        printf("Reste      : %d %% %d = %d\n", a, b, a % b);
    } else {
        printf("Erreur : division par zero impossible.\n");
    }

    return 0;
}
\end{lstlisting}

\textbf{Points importants :}
\begin{itemize}[nosep]
  \item En C, la division entière \texttt{a/b} tronque vers zéro.
  \item L'opérateur \texttt{\%} donne le reste de la division euclidienne.
  \item Il faut vérifier que $b\neq0$ avant de diviser.
  \item \texttt{scanf} prend l'adresse des variables avec \texttt{\&}.
  \item Pour afficher le caractère \texttt{\%} avec \texttt{printf}, on écrit \texttt{\%\%}.
\end{itemize}
\end{boitecorrige}

\subsection*{Exercice 2 - Structures conditionnelles \niveau{\facile}}

\begin{boiteexercice}
Écrire un programme C qui :
\begin{enumerate}
  \item Lit un entier $n$.
  \item Affiche si $n$ est positif, négatif ou nul.
  \item Détermine si $n$ est pair ou impair.
  \item Affiche la valeur absolue de $n$.
\end{enumerate}
\end{boiteexercice}

\begin{boitecorrige}
\begin{lstlisting}
#include <stdio.h>

int main() {
    int n;

    printf("Entrez un entier : ");
    scanf("%d", &n);

    /* Signe */
    if (n > 0) {
        printf("%d est positif.\n", n);
    } else if (n < 0) {
        printf("%d est negatif.\n", n);
    } else {
        printf("Le nombre est nul.\n");
    }

    /* Parite */
    if (n % 2 == 0) {
        printf("%d est pair.\n", n);
    } else {
        printf("%d est impair.\n", n);
    }

    /* Valeur absolue */
    int abs_n = (n >= 0) ? n : -n;
    printf("Valeur absolue de %d : %d\n", n, abs_n);

    return 0;
}
\end{lstlisting}

L'opérateur ternaire \texttt{(cond) ? val\_vrai : val\_faux} est une façon concise d'écrire un \texttt{if-else} sur une expression.
\end{boitecorrige}

\section{Boucles et tableaux}

\subsection*{Exercice 3 - Boucles et calculs \niveau{\moyen}}

\begin{boiteexercice}
Écrire des programmes C pour :
\begin{enumerate}[label=(\alph*)]
  \item Calculer la factorielle $n!$ pour un $n$ saisi par l'utilisateur.
  \item Afficher la table de multiplication de 1 à 10 pour un entier $k$ donné.
  \item Calculer la somme $S = 1 + 2 + \cdots + n$ et vérifier avec la formule $n(n+1)/2$.
  \item Déterminer si un entier $n$ est premier.
\end{enumerate}
\end{boiteexercice}

\begin{boitecorrige}
\textbf{(a) Factorielle :}
\begin{lstlisting}
#include <stdio.h>
int main() {
    int n;
    long long fact = 1;
    printf("n = "); scanf("%d", &n);
    for (int i = 1; i <= n; i++) {
        fact *= i;
    }
    printf("%d! = %lld\n", n, fact);
    return 0;
}
\end{lstlisting}

\textbf{(b) Table de multiplication :}
\begin{lstlisting}
int k;
printf("k = "); scanf("%d", &k);
for (int i = 1; i <= 10; i++) {
    printf("%d x %d = %d\n", k, i, k * i);
}
\end{lstlisting}

\textbf{(c) Somme avec vérification :}
\begin{lstlisting}
int n, S = 0;
printf("n = "); scanf("%d", &n);
for (int i = 1; i <= n; i++) S += i;
printf("Somme boucle   : %d\n", S);
printf("Formule n(n+1)/2 : %d\n", n * (n + 1) / 2);
\end{lstlisting}

\textbf{(d) Test de primalité :}
\begin{lstlisting}
#include <stdio.h>
int est_premier(int n) {
    if (n < 2) return 0;
    for (int i = 2; i * i <= n; i++) {
        if (n % i == 0) return 0;  /* divisible */
    }
    return 1;
}
int main() {
    int n;
    printf("n = "); scanf("%d", &n);
    if (est_premier(n))
        printf("%d est premier.\n", n);
    else
        printf("%d n'est pas premier.\n", n);
    return 0;
}
\end{lstlisting}
L'optimisation clé : on teste les diviseurs jusqu'à $\sqrt{n}$ seulement.
\end{boitecorrige}

\subsection*{Exercice 4 - Tableaux et tri \niveau{\moyen}}

\begin{boiteexercice}
\begin{enumerate}
  \item Écrire un programme qui remplit un tableau de $n$ entiers, puis calcule et affiche la somme, la moyenne, le maximum et le minimum.
  \item Implémenter le tri à bulles (\textit{bubble sort}) pour trier un tableau dans l'ordre croissant.
\end{enumerate}
\end{boiteexercice}

\begin{boitecorrige}
\textbf{1.}
\begin{lstlisting}
#include <stdio.h>
#define N 10
int main() {
    int tab[N], somme = 0, max, min;
    printf("Entrez %d entiers :\n", N);
    for (int i = 0; i < N; i++) scanf("%d", &tab[i]);

    max = min = tab[0];
    for (int i = 0; i < N; i++) {
        somme += tab[i];
        if (tab[i] > max) max = tab[i];
        if (tab[i] < min) min = tab[i];
    }
    printf("Somme   = %d\n", somme);
    printf("Moyenne = %.2f\n", (double)somme / N);
    printf("Max     = %d\n", max);
    printf("Min     = %d\n", min);
    return 0;
}
\end{lstlisting}

\textbf{2. Tri à bulles :}
\begin{lstlisting}
void tri_bulles(int tab[], int n) {
    int tmp;
    for (int i = 0; i < n - 1; i++) {
        for (int j = 0; j < n - 1 - i; j++) {
            if (tab[j] > tab[j + 1]) {
                /* echange */
                tmp = tab[j];
                tab[j] = tab[j + 1];
                tab[j + 1] = tmp;
            }
        }
    }
}
\end{lstlisting}

Le tri à bulles a une complexité $O(n^2)$ : à chaque passage, le plus grand élément non trié remonte à sa place (comme une bulle).
\end{boitecorrige}

\section{Fonctions et pointeurs}

\subsection*{Exercice 5 - Fonctions \niveau{\moyen}}

\begin{boiteexercice}
\begin{enumerate}
  \item Écrire une fonction \texttt{puissance(double base, int n)} qui calcule $\text{base}^n$ par multiplications successives.
  \item Écrire une fonction récursive \texttt{fibonacci(int n)} qui calcule le $n$-ième terme de la suite de Fibonacci définie par $F_0=0$, $F_1=1$, $F_n=F_{n-1}+F_{n-2}$.
  \item Écrire une fonction \texttt{pgcd(int a, int b)} qui calcule le PGCD par l'algorithme d'Euclide.
\end{enumerate}
\end{boiteexercice}

\begin{boitecorrige}
\begin{lstlisting}
#include <stdio.h>

/* (1) Puissance iterative */
double puissance(double base, int n) {
    double resultat = 1.0;
    int exp = (n >= 0) ? n : -n;
    for (int i = 0; i < exp; i++) resultat *= base;
    return (n >= 0) ? resultat : 1.0 / resultat;
}

/* (2) Fibonacci recursif */
long long fibonacci(int n) {
    if (n == 0) return 0;
    if (n == 1) return 1;
    return fibonacci(n - 1) + fibonacci(n - 2);
}

/* (3) PGCD par algorithme d'Euclide */
int pgcd(int a, int b) {
    while (b != 0) {
        int r = a % b;
        a = b;
        b = r;
    }
    return a;
}

int main() {
    printf("2^10 = %.0f\n", puissance(2, 10));
    printf("F(10) = %lld\n", fibonacci(10));
    printf("pgcd(48, 18) = %d\n", pgcd(48, 18));
    return 0;
}
\end{lstlisting}

\textbf{Attention :} la version récursive de Fibonacci est très inefficace pour les grands $n$ (complexité exponentielle). En pratique, on utilise une version itérative ou la mémoïsation.
\end{boitecorrige}

\subsection*{Exercice 6 - Pointeurs et passage par adresse \niveau{\difficile}}

\begin{boiteexercice}
\begin{enumerate}
  \item Écrire une fonction \texttt{echange(int *a, int *b)} qui échange les valeurs de deux variables. Expliquer pourquoi le passage par valeur ne fonctionnerait pas.
  \item Écrire une fonction qui calcule simultanément le minimum et le maximum d'un tableau et les retourne via des pointeurs.
  \item Qu'affiche le programme suivant ?
\begin{lstlisting}
int main() {
    int t[] = {10, 20, 30, 40, 50};
    int *p = t;
    printf("%d\n", *p);
    printf("%d\n", *(p + 2));
    p++;
    printf("%d\n", *p);
    printf("%d\n", p[2]);
    return 0;
}
\end{lstlisting}
\end{enumerate}
\end{boiteexercice}

\begin{boitecorrige}
\textbf{1.}
\begin{lstlisting}
void echange(int *a, int *b) {
    int tmp = *a;
    *a = *b;
    *b = tmp;
}
\end{lstlisting}
Si on passait par valeur (\texttt{void echange(int a, int b)}), la fonction travaillerait sur des copies locales ; les variables originales dans \texttt{main} resteraient inchangées. Avec les pointeurs, on accède directement aux emplacements mémoire d'origine.

\textbf{2.}
\begin{lstlisting}
void min_max(int tab[], int n, int *min, int *max) {
    *min = *max = tab[0];
    for (int i = 1; i < n; i++) {
        if (tab[i] < *min) *min = tab[i];
        if (tab[i] > *max) *max = tab[i];
    }
}
/* Appel : int mn, mx; min_max(t, 5, &mn, &mx); */
\end{lstlisting}

\textbf{3.} Analyse pas à pas :
\begin{itemize}[nosep]
  \item \texttt{p = t} : $p$ pointe sur \texttt{t[0]} (valeur 10).
  \item \texttt{*p} : valeur à l'adresse $p$ = \textbf{10}.
  \item \texttt{*(p+2)} : adresse de \texttt{t[2]}, valeur = \textbf{30}.
  \item \texttt{p++} : $p$ avance d'un entier, pointe maintenant sur \texttt{t[1]} (valeur 20).
  \item \texttt{*p} = \textbf{20}.
  \item \texttt{p[2]} $\equiv$ \texttt{*(p+2)} : deux cases après \texttt{t[1]}, soit \texttt{t[3]} = \textbf{40}.
\end{itemize}
Le programme affiche : 10, 30, 20, 40.
\end{boitecorrige}


\subsection*{Exercice 7 - Récursivité : suite de Fibonacci \niveau{\moyen}}

\begin{boiteexercice}
La suite de Fibonacci est définie par $F_0=0$, $F_1=1$, $F_n=F_{n-1}+F_{n-2}$ pour $n\geq2$.
\begin{enumerate}
  \item Écrire une fonction C récursive \texttt{fib\_rec(int n)} qui calcule $F_n$.
  \item Écrire une fonction C itérative \texttt{fib\_iter(int n)} qui calcule $F_n$.
  \item Comparer les complexités temporelles des deux approches. Pourquoi la version récursive est-elle très lente pour $n$ grand ?
  \item Écrire un programme \texttt{main} qui affiche les 15 premiers termes avec les deux méthodes.
  \item Modifier \texttt{fib\_iter} pour qu'elle calcule également la somme $S_n = \sum_{k=0}^n F_k$.
\end{enumerate}
\end{boiteexercice}

\begin{boitecorrige}
\textbf{1.} Version récursive :
\begin{lstlisting}
int fib_rec(int n) {
    if (n <= 0) return 0;
    if (n == 1) return 1;
    return fib_rec(n-1) + fib_rec(n-2);
}
\end{lstlisting}

\textbf{2.} Version itérative :
\begin{lstlisting}
int fib_iter(int n) {
    if (n <= 0) return 0;
    if (n == 1) return 1;
    int a = 0, b = 1, c;
    for (int i = 2; i <= n; i++) {
        c = a + b;
        a = b;
        b = c;
    }
    return b;
}
\end{lstlisting}

\textbf{3.} Complexité :
\begin{itemize}[nosep]
  \item \texttt{fib\_rec} : complexité $O(\varphi^n)$ où $\varphi=(1+\sqrt{5})/2\approx1{,}618$ (nombre d'or). Chaque appel engendre 2 sous-appels ; de nombreux sous-problèmes sont recalculés plusieurs fois (arbre d'appels exponentiel).
  \item \texttt{fib\_iter} : complexité $O(n)$ en temps, $O(1)$ en espace. Elle est donc exponentiellement plus rapide.
\end{itemize}
Pour $n=40$, \texttt{fib\_rec} effectue $\approx 10^8$ appels tandis que \texttt{fib\_iter} effectue 40 itérations.

\textbf{4.} Programme principal :
\begin{lstlisting}
#include <stdio.h>
int main(void) {
    printf("n   fib_rec  fib_iter\n");
    for (int i = 0; i < 15; i++)
        printf("%2d  %6d   %6d\n", i, fib_rec(i), fib_iter(i));
    return 0;
}
\end{lstlisting}

\textbf{5.} Calcul simultané de la somme ($S_n = F_{n+2}-1$) :
\begin{lstlisting}
int fib_iter_sum(int n, int *somme) {
    if (n <= 0) { *somme = 0; return 0; }
    int a = 0, b = 1, c, s = 1;
    for (int i = 2; i <= n; i++) {
        c = a + b; a = b; b = c; s += c;
    }
    *somme = s;
    return b;
}
\end{lstlisting}
\end{boitecorrige}

\subsection*{Exercice 8 - Tableaux 2D et algèbre matricielle \niveau{\moyen}}

\begin{boiteexercice}
On travaille avec des matrices $4\times4$ stockées comme tableaux 2D en C.
\begin{enumerate}
  \item Écrire une fonction \texttt{init\_identite(int M[4][4])} qui initialise $M$ à la matrice identité.
  \item Écrire une fonction \texttt{trace(int M[4][4])} qui retourne la trace de $M$.
  \item Écrire une fonction \texttt{transposee(int A[4][4], int B[4][4])} qui stocke dans $B$ la transposée de $A$.
  \item Écrire une fonction \texttt{produit\_mat\_vec(int A[4][4], int v[4], int res[4])} qui calcule $\vec{res} = A\vec{v}$.
  \item Tester avec la matrice $A_{ij}=i+j$ ($i,j$ de 0 à 3) et le vecteur $\vec{v}=(1,0,0,0)^T$.
\end{enumerate}
\end{boiteexercice}

\begin{boitecorrige}
\textbf{1.}
\begin{lstlisting}
void init_identite(int M[4][4]) {
    for (int i = 0; i < 4; i++)
        for (int j = 0; j < 4; j++)
            M[i][j] = (i == j) ? 1 : 0;
}
\end{lstlisting}

\textbf{2.}
\begin{lstlisting}
int trace(int M[4][4]) {
    int t = 0;
    for (int i = 0; i < 4; i++) t += M[i][i];
    return t;
}
\end{lstlisting}

\textbf{3.}
\begin{lstlisting}
void transposee(int A[4][4], int B[4][4]) {
    for (int i = 0; i < 4; i++)
        for (int j = 0; j < 4; j++)
            B[j][i] = A[i][j];
}
\end{lstlisting}

\textbf{4.}
\begin{lstlisting}
void produit_mat_vec(int A[4][4], int v[4], int res[4]) {
    for (int i = 0; i < 4; i++) {
        res[i] = 0;
        for (int j = 0; j < 4; j++)
            res[i] += A[i][j] * v[j];
    }
}
\end{lstlisting}

\textbf{5.} Pour $A_{ij}=i+j$ : $A = \begin{pmatrix}0&1&2&3\\1&2&3&4\\2&3&4&5\\3&4&5&6\end{pmatrix}$.
Trace : $0+2+4+6=12$. Produit $A\vec{v}$ avec $\vec{v}=(1,0,0,0)^T$ donne la première colonne de $A$ : $(0,1,2,3)^T$.
\begin{lstlisting}
int main(void) {
    int A[4][4], v[4]={1,0,0,0}, res[4], B[4][4];
    for (int i=0;i<4;i++) for(int j=0;j<4;j++) A[i][j]=i+j;
    printf("Trace = %d\n", trace(A));          /* 12 */
    produit_mat_vec(A, v, res);
    printf("Av = (%d,%d,%d,%d)\n", res[0],res[1],res[2],res[3]);
    return 0;
}
\end{lstlisting}
\end{boitecorrige}

\subsection*{Exercice 9 - Allocation dynamique et pointeurs sur fonctions \niveau{\difficile}}

\begin{boiteexercice}
\begin{enumerate}
  \item Écrire une fonction \texttt{swap(int *a, int *b)} qui échange les valeurs de deux variables entières. Expliquer pourquoi le passage par pointeur est nécessaire.
  \item Écrire une fonction \texttt{creer\_tableau(int n)} qui alloue dynamiquement un tableau de $n$ entiers, les initialise à $0,1,2,\ldots,n-1$ et retourne le pointeur.
  \item Écrire une fonction \texttt{appliquer(int *t, int n, int (*f)(int))} qui applique la fonction \texttt{f} à chaque élément du tableau.
  \item Tester avec les fonctions $f(x)=x^2$ et $g(x)=2x+1$ sur un tableau de 6 éléments.
  \item Expliquer ce que fait le code suivant et identifier l'erreur :\\
  \texttt{int *p = malloc(10*sizeof(int)); p[10] = 5;}
\end{enumerate}
\end{boiteexercice}

\begin{boitecorrige}
\textbf{1.}
\begin{lstlisting}
void swap(int *a, int *b) {
    int tmp = *a;
    *a = *b;
    *b = tmp;
}
\end{lstlisting}
Sans pointeurs, C passe les arguments par valeur : la fonction travaillerait sur des copies et les variables originales resteraient inchangées.

\textbf{2.}
\begin{lstlisting}
int *creer_tableau(int n) {
    int *t = malloc(n * sizeof(int));
    if (t == NULL) { fprintf(stderr, "Erreur malloc\n"); exit(1); }
    for (int i = 0; i < n; i++) t[i] = i;
    return t;
}
\end{lstlisting}
Remarque : le tableau alloué doit être libéré par l'appelant avec \texttt{free(t)}.

\textbf{3.}
\begin{lstlisting}
void appliquer(int *t, int n, int (*f)(int)) {
    for (int i = 0; i < n; i++) t[i] = f(t[i]);
}
\end{lstlisting}

\textbf{4.}
\begin{lstlisting}
int carre(int x) { return x * x; }
int double_plus_un(int x) { return 2*x + 1; }

int main(void) {
    int *t = creer_tableau(6);  /* {0,1,2,3,4,5} */
    appliquer(t, 6, carre);     /* {0,1,4,9,16,25} */
    for (int i=0;i<6;i++) printf("%d ", t[i]);
    free(t);
    t = creer_tableau(6);
    appliquer(t, 6, double_plus_un); /* {1,3,5,7,9,11} */
    for (int i=0;i<6;i++) printf("%d ", t[i]);
    free(t);
    return 0;
}
\end{lstlisting}

\textbf{5.} \texttt{malloc(10*sizeof(int))} alloue 10 entiers aux indices $0$ à $9$. \texttt{p[10]=5} écrit en dehors de la zone allouée (\textbf{dépassement de tampon} -- \textit{buffer overflow}). C'est un comportement indéfini qui peut provoquer une corruption mémoire, un segfault ou des failles de sécurité.
\end{boitecorrige}

\subsection*{Exercice 10 - Lecture de fichier et statistiques \niveau{\moyen}}

\begin{boiteexercice}
Un fichier texte \texttt{donnees.txt} contient des entiers séparés par des espaces ou sauts de ligne.
\begin{enumerate}
  \item Écrire une fonction \texttt{lire\_fichier(const char *nom, int *t, int max)} qui lit les entiers du fichier dans le tableau \texttt{t} (de taille maximale \texttt{max}) et retourne le nombre de valeurs lues.
  \item Écrire une fonction \texttt{moyenne(int *t, int n)} qui retourne la moyenne sous forme de \texttt{double}.
  \item Écrire une fonction \texttt{ecart\_type(int *t, int n, double moy)} qui retourne l'écart-type $\sigma = \sqrt{\frac{1}{n}\sum_{i=1}^n(x_i-\bar x)^2}$.
  \item Écrire le programme principal qui affiche $n$, $\bar x$, $\sigma_x$, le minimum et le maximum.
\end{enumerate}
\end{boiteexercice}

\begin{boitecorrige}
\textbf{1.}
\begin{lstlisting}
#include <stdio.h>
#include <math.h>

int lire_fichier(const char *nom, int *t, int max) {
    FILE *f = fopen(nom, "r");
    if (f == NULL) { perror("fopen"); return -1; }
    int n = 0;
    while (n < max && fscanf(f, "%d", &t[n]) == 1) n++;
    fclose(f);
    return n;
}
\end{lstlisting}

\textbf{2.}
\begin{lstlisting}
double moyenne(int *t, int n) {
    double s = 0.0;
    for (int i = 0; i < n; i++) s += t[i];
    return s / n;
}
\end{lstlisting}

\textbf{3.}
\begin{lstlisting}
double ecart_type(int *t, int n, double moy) {
    double s = 0.0;
    for (int i = 0; i < n; i++) {
        double d = t[i] - moy;
        s += d * d;
    }
    return sqrt(s / n);
}
\end{lstlisting}

\textbf{4.}
\begin{lstlisting}
int main(void) {
    int t[1000];
    int n = lire_fichier("donnees.txt", t, 1000);
    if (n <= 0) return 1;
    double moy = moyenne(t, n);
    double sigma = ecart_type(t, n, moy);
    int mn = t[0], mx = t[0];
    for (int i = 1; i < n; i++) {
        if (t[i] < mn) mn = t[i];
        if (t[i] > mx) mx = t[i];
    }
    printf("n     = %d\n", n);
    printf("moy   = %.4f\n", moy);
    printf("sigma = %.4f\n", sigma);
    printf("min   = %d, max = %d\n", mn, mx);
    return 0;
}
\end{lstlisting}
\end{boitecorrige}

\subsection*{Exercice 11 - Structures et tri de points \niveau{\moyen}}

\begin{boiteexercice}
On souhaite manipuler des points du plan.
\begin{enumerate}
  \item Définir une structure \texttt{Point} avec deux champs \texttt{double x} et \texttt{double y}.
  \item Écrire une fonction \texttt{distance(Point a, Point b)} qui retourne la distance euclidienne entre deux points.
  \item Écrire une fonction de comparaison \texttt{cmp\_x} utilisable avec \texttt{qsort} pour trier un tableau de points par abscisse croissante.
  \item Dans \texttt{main}, créer un tableau de 5 points, le trier par abscisse avec \texttt{qsort}, puis afficher les points triés.
  \item Écrire une fonction \texttt{barycentre(Point *t, int n)} qui retourne le barycentre de $n$ points.
\end{enumerate}
\end{boiteexercice}

\begin{boitecorrige}
\textbf{1.}
\begin{lstlisting}
#include <stdio.h>
#include <stdlib.h>
#include <math.h>

typedef struct {
    double x, y;
} Point;
\end{lstlisting}

\textbf{2.}
\begin{lstlisting}
double distance(Point a, Point b) {
    double dx = a.x - b.x;
    double dy = a.y - b.y;
    return sqrt(dx*dx + dy*dy);
}
\end{lstlisting}

\textbf{3.} La fonction de comparaison pour \texttt{qsort} reçoit des pointeurs génériques \texttt{void*} :
\begin{lstlisting}
int cmp_x(const void *p1, const void *p2) {
    const Point *a = (const Point *)p1;
    const Point *b = (const Point *)p2;
    if (a->x < b->x) return -1;
    if (a->x > b->x) return  1;
    return 0;
}
\end{lstlisting}

\textbf{4.}
\begin{lstlisting}
int main(void) {
    Point pts[5] = {{3,1},{1,4},{4,1},{1,5},{9,2}};
    qsort(pts, 5, sizeof(Point), cmp_x);
    for (int i = 0; i < 5; i++)
        printf("(%.1f, %.1f)\n", pts[i].x, pts[i].y);
    /* Affiche dans l'ordre des x : 1,1,3,4,9 */
    return 0;
}
\end{lstlisting}

\textbf{5.}
\begin{lstlisting}
Point barycentre(Point *t, int n) {
    Point g = {0.0, 0.0};
    for (int i = 0; i < n; i++) {
        g.x += t[i].x;
        g.y += t[i].y;
    }
    g.x /= n;
    g.y /= n;
    return g;
}
\end{lstlisting}
Pour les 5 points ci-dessus, le barycentre est $G=(18/5, 13/5)=(3{,}6;\;2{,}6)$.
\end{boitecorrige}

\ifdefined\inMainDoc\else\end{document}\fi
